\section{问题重述}

\subsection{问题背景}

无线电干扰源的定位与清除是频谱管理的核心任务，通常经历初步监测、区域性排查与精准定位三个阶段。前两个阶段可由固定监测站、监测车或无人机完成，可把干扰源圈定到一片区域；但受环境遮挡与低功率辐射的限制，最后的精准定位仍须由技术人员携带便携式测向机徒步完成，效率低且在危险环境中风险大。本题考察的正是这最后一环：用搭载测向机、光学探测仪与激光枪的机器狗，在一个已知的圆形目标区域内自动完成\textbf{发现—定位—清除}的全过程，并使总耗时尽可能短。

目标区域是以原点为心、半径 $R=1800\mmet$ 的圆盘，正东为 $x$ 轴正向、正北为 $y$ 轴正向。区域中干扰源个数未知（第三、四问为 $10\sim16$ 个），每个源占用 $1\sim20$ 中互不相同且不随时间变化的一个频道，其位置、有效接收半径（介于 $1000\mmet$ 与 $1500\mmet$ 之间）以及（第四问的）定向方向都不可直接读取。机器狗只能通过 \texttt{/measure} 与 \texttt{/clear} 两条指令与环境交互，这使问题本质上是一个\textbf{部分可观测的在线决策问题}\cite{kaelbling}，而不是已知坐标的路径规划问题。

\subsection{需要解决的问题}

\begin{itemize}
  \item \textbf{问题一（定位区域的直径与覆盖判定）.} 给定若干检测点及其对同一干扰源的示向度，要求给出计算多边形定位区域\emph{直径}的算法，并判定以该直径为直径的圆能否覆盖该区域。其数学本质是\emph{角度误差约束下的凸多边形度量问题}。

  \item \textbf{问题二（第二检测点的选择）.} 已知某\emph{全向}源在首个检测点处的示向度，要求给出第二个检测点的选择策略及其可行的\emph{候选区域}。其本质是在“第二点必能收到信号”这一\emph{对未知量一致成立}的约束下，权衡定位精度与移动代价。

  \item \textbf{问题三（全向源的搜索、定位与清除）.} 区域内共有 $10\sim16$ 个全向源而真实总数未知，机器狗须在\emph{保证全部清除}的前提下使任务时间最短。其本质是一个\emph{带覆盖保证的部分可观测在线搜索与调度问题}：核心在于用最小测站集证明搜索完备性，并以在线贪心策略压缩固定与移动成本。

  \item \textbf{问题四（混合全向源与定向源）.} 源中同时含有定向源，其定向方向未知、有效覆盖为方向两侧各 $90^\circ$，其余条件同问题三。其本质是把问题三的\emph{距离覆盖}升级为\emph{位置与朝向联合的方向覆盖}。
\end{itemize}

\subsection{关键物理与计时规则}

题目与附件给出的规则构成了模型的全部硬约束，汇总于\cref{tab:rules}。其中有三条特别影响建模：

\begin{enumerate}
  \item \textbf{测向误差是有界的、确定性的，而不是随机的.} 题目只保证误差落在 $[-1^\circ,1^\circ]$ 内，且“同一地点重复检测不会改变检测误差”。因此不能假设零均值、独立或高斯分布，也不应使用最小二乘或卡尔曼滤波这类依赖概率模型的估计器；正确的处理是\textbf{集员（set-membership）定位}\cite{milanese}。
  \item \textbf{清除的判据是距离而不是方向.} 光学精确定位与激光清除只要求清除点与源的距离不超过 $20\mmet$，与信号覆盖角度范围无关。这使得“即使再也收不到某定向源的信号，只要还持有它的一次方位约束，仍能用光学覆盖把它清掉”成为可能，是第四问兜底机制的基础。
  \item \textbf{移动、切频与检测的耗时结构固定.} 移动 $5\mmet/\msec$、切频 $1\msec$、检测 $5\msec$、清除失败 $3\msec$、清除成功 $5\msec$。这些常数决定了后文的目标函数与成本权衡：一次额外检测（$5\sim6\msec$）约相当于 $25\sim30\mmet$ 的移动，因此顺路增加一次检测通常是划算的。
\end{enumerate}

\begin{table}[!htbp]
  \caption{题目给定的物理规则与计时常数}\label{tab:rules}\centering
  \begin{tabular}{llc}
    \toprule[1.5pt]
    项目 & 规则 & 取值 \\
    \midrule[1pt]
    目标区域    & 以原点为心的圆盘，东为 $x$ 正向、北为 $y$ 正向 & $R=1800\mmet$ \\
    源个数      & 第三、四问，不通过接口返回                     & $10\sim16$ \\
    频道        & 整数，一源一频道，互不相同且不变               & $1\sim20$ \\
    示向度误差  & 全局有界；同点重复检测误差固定                 & $\pm1^\circ$ \\
    有效接收半径& 每源一值，未知                                 & $1000\sim1500\mmet$ \\
    定向覆盖    & 定向方向两侧各 $90^\circ$，含边界（闭半平面）  & $180^\circ$ \\
    移动        & 直线，可走出目标区域                           & $5\mmet/\msec$ \\
    检测        & 停止移动后调整天线至稳定读数                   & $5\msec$ \\
    切换频道    & 任意两频道之间                                 & $1\msec$ \\
    光学清除    & 距离 $\le20\mmet$ 必成功，与辐射方向无关       & 成功 $5\msec$／失败 $3\msec$ \\
    近场        & 距离 $\le5\mmet$ 且在覆盖角内，返回 \texttt{near} & $5\mmet$ \\
    虚拟时间上限& 防止死循环                                     & $360000\msec$ \\
    程序运行时间& $25$ 分钟窗口与 $20$ 分钟程序时限取早者         & $\le1200\msec$ \\
    \bottomrule[1.5pt]
  \end{tabular}
\end{table}
